\section{Conclusion}
\label{sec:conclusion}

\subsection{What the investigation establishes}

The starting observation was that an optimiser can succeed while a registration
fails. The investigation traced this to one structural fact: the scan supplies
surface geometry, whereas anatomical correspondence is an intended consequence
of the recovered parameters, never an observed quantity. Each result is a
different consequence of that fact. Input quality is not the limit (\F{F2}).
Surface proximity is an unreliable proxy for anatomy, because free deformation
lets a wrong explanation hide behind a good surface (\F{F1}).
Correspondence splits into identity and localisation (\F{F3}); initialisation
helps where pose alone determines placement and not where attachment depends on
shape (\F{F4}); thin structures are under-weighted before any pose question
arises (\F{F5}); and the objective has no term observing joint location, so
supervision acts only where it is supplied (\F{F6}). The residual error is
structured by joint and specimen, not by chance (\F{F8}), and the reference
used to see all of this had itself to be audited (\F{F7}).

Four distinctions summarise the outcome: representability is not
recoverability; regional identity is not within-region location; plausibility
is not correspondence; and fit quality is not measurement validity. Each was the
concrete reason that a plausible intervention failed. The contribution is
therefore not a fitting trick but a technical framework that makes articulated
3D reconstruction measurable, diagnostically interpretable and independently
verifiable. A useful parametric representation is defined not only by what it
can represent but by what anatomical meaning survives its recovery from data:
a reconstruction is not yet a phenotype.

\subsection{Relation to existing knowledge}

Several findings confirm established principles in a new domain and are not
claimed as discoveries: the dissociation between fit error and target error
(known from medical registration~\cite{fitzpatrick}), the inability of
intrinsic descriptors to break mirror symmetry~\cite{donati}, the expectation
that dense-correspondence models become multi-modal on symmetric
surfaces~\cite{haugaard,cse}, and the gap between benchmark success and system
success~\cite{liao}. What this work adds is their measurement, under
pre-registered single-factor ablation and an independently audited reference,
in a deformable articulated animal model, together with per-region numbers that
decide which traits can be trusted.

\subsection{Scope of the evidence}

The conclusions are bounded by the data in a precise way. The expert benchmark
has 11 specimens, so the pre-registered specimen-level test resolves effects
that are consistent across about ten of eleven specimens, and smaller effects
are carried by intervals and the joint-level model. Synthetic targets come from
the model and therefore give a lower bound for real scans. A fixed topology and a
25-dimensional shape basis define what the model can express, and residual
deformation is ambiguous with pose by construction (Section~\ref{sec:formal}).
Annotation has its own repeatability, trait definitions are coupled to scale,
and a population-level result does not carry over automatically to individual
specimens. The ground-truth audit retired earlier analyses built on the
superseded reference; every result reported here rests on the corrected
benchmark or on exact synthetic ground truth.

\subsection{Outlook}

Each direction follows from a diagnosed failure, and none was implemented.
Independent pointwise votes collapse, which calls for a global mechanism such
as functional maps~\cite{fmaps} or optimal transport~\cite{ot}, with partial
variants for incomplete scans. The absence of an anatomical signal in the
objective calls for skeleton-conditioned correspondence or a learned 3D joint
predictor; deformation that absorbs pose error calls for structured deformation;
and the synthetic-to-real gap calls for scan-realistic synthetic training. The
most ambitious direction is an anatomical canonical field that predicts, for
every observed point, its part, joint-relative coordinates, laterality and a
confidence, so that correspondence becomes the mapping of a scan into a
canonical anatomical coordinate system: exactly the information the diagnosed
objective lacks. The failure mechanisms are not specific to ants; whether they
recur on other body plans is an empirical question that the same framework is
built to answer.

\subsection{What the internship taught me}

Technically, it taught me to treat an objective as a proxy and to test the
mechanism behind it: to build ground truth, to audit it before trusting it, and
to report negative results with the same care as positive ones. Professionally,
it taught me to work on a live shared repository under review, to fix success
criteria before an experiment, to run reproducible experiments on two HPC
systems, and to redirect a project when the evidence contradicts its starting
hypothesis, a decision that is only defensible alongside the evidence that
motivated it.
